\section*{Hoofdstuk \ref{ch:b3:photons-phonons} --- Fotonen en fononen}

\begin{solution}{exo:b3:photons-phonons:1}
(a) Uit $\dd u/\dd\lambda = 0$ volgt $x\eu^x/(\eu^x - 1) = 5$,
$x = hc/\lambda k_{\text{B}}T = 4.965$: $\lambda_{\max}T =
hc/4.965k_{\text{B}} = \qty{2898}{\micro m.K}$. (b) $0.50$,
$1.07$, $9.35\,\unit{\micro m}$ en \qty{1.06}{mm}. (c) Alleen die
van de Zon; een gloeidraad piekt in het nabije infrarood en levert
maar een paar procent van zijn vermogen als licht --- de
inefficiëntie die hem uit de handel nam. (d) Mensen en gebouwen
``stralen'' bij tien micrometer: warmtebeeldcamera's zijn camera's
voor de tweede piek.
\end{solution}

\begin{solution}{exo:b3:photons-phonons:2}
(a) $\sigma A(T_1^4 - T_2^4) \approx \qty{135}{W}$. (b) De
retourflux van de kamer zit al in het verschil verwerkt; kleding
schuift warme tussenoppervlakken in --- straling is een grootboek
met twee zijden. (c)
$T_{\text{eff}} = [(0.7 \times 1361)/(4\sigma)]^{1/4} =
\qty{255}{K}$. (d) De bonus van \qty{33}{K} is het broeikaseffect:
de infrarode deken van
\cref{pb:b3:harmonic-oscillator:1}, gekwantificeerd in
\cref{exo:b3:photons-phonons:9}.
\end{solution}

\begin{solution}{exo:b3:photons-phonons:3}
(a) $\num{2.03e7} \times 2.725^3 \approx \qty{4.1e8}{m^{-3}} =
411$ per cm$^3$. (b) $1361/\qty{2.2e-19}{J} \approx
\num{6e21}$ fotonen per vierkante meter per seconde. (c) $3\,\unit{W}/
\qty{3.6e-19}{J} \approx \num{8e18}$ per seconde. (d) Astronomische
signalen komen foton voor foton binnen --- tempo's die een teller
kan onderscheiden; de $10^{22}$ per seconde van een kachel zijn
voor elke detector een continuüm.
\end{solution}

\begin{solution}{exo:b3:photons-phonons:4}
(a) $k_{\text{B}}T$ per mode maal het aantal moden $\nu^2$:
geïntegreerd oneindig --- de catastrofe. (b) Voor $h\nu \ll
k_{\text{B}}T$. (c) $h\nu/k_{\text{B}}T = 0.025$: diep in de
klassieke staart, waar de intensiteit $\propto T$ is --- vandaar
de ``helderheidstemperatuur'' in kelvin. (d) $2.8$: bij de piek
werkt Plancks kwantumkorting op volle kracht.
\end{solution}

\begin{solution}{exo:b3:photons-phonons:5}
(a) Voor de gerichte zonnebundel is $P \approx \Phi/c =
\qty{4.5}{\micro Pa}$: elf ordegrootten onder de atmosfeer ---
zeilen werken alleen doordat de ruimte al het andere aftrekt. (b)
$aT^4/3 \approx \qty{1.3e13}{Pa}$: minder dan één duizendste
van de gasdruk --- de Zon wordt door gas gedragen. (c)
$P_{\text{rad}}/P_{\text{gas}} \propto T^3/n$: zwaardere sterren
branden heter, en de $T^4$-vloed gaat uiteindelijk wedijveren met
de greep van de zwaartekracht op het gas --- lichtsterke
instabiliteit begrenst sterren rond
$\sim100\,M_\odot$. (d) $F = P/c = \qty{3.3}{\micro N}$:
trekstralen blijven fictie; pincetten gebruiken veld\emph{gradiënten}
op microscopische objecten, waar piconewtons regeren
(\cref{pb:b3:microcanonical-ensemble:1}).
\end{solution}

\begin{solution}{exo:b3:photons-phonons:6}
(a) $1 + z = 3000/2.725 \approx 1100$. (b) De $\lambda$ van elke mode
rekt met dezelfde factor: de vorm $\nu^3$ beeldt op zichzelf af
met $T \to T/(1+z)$ --- een planckkromme laat zich niet
ont-plancken door uitdijing. (c) De diepdode piek bij
$\approx\qty{0.97}{\micro m}$ komt aan bij \qty{1.06}{mm}. (d) Dezelfde
hemel is leeg bij \qty{0.5}{\micro m} en verzadigd bij een
millimeter: duisternis is een uitspraak over golflengte.
\end{solution}

\begin{solution}{exo:b3:photons-phonons:7}
(a) $\theta_{\text{D}} \approx \qty{500}{K}$ tegenover de gefitte
\qty{343}{K}: de juiste schaal (een net gemiddelde over
longitudinale en transversale snelheden dicht het gat). (b)
$\approx\qty{2200}{K}$ --- kloppend: diamant is het uiterste
geval. (c) Een lage geluidssnelheid (zacht) en zware atomen
drukken $\theta_{\text{D}}$ allebei omlaag: het rooster van lood
is vrijwel meteen klassiek. (d) $\theta_{\text{D}} \propto
c_{\text{s}}n^{1/3} \propto \sqrt{k/m}$-achtig: stijf en licht
bevriest laat, zacht en zwaar bevriest vroeg.
\end{solution}

\begin{solution}{exo:b3:photons-phonons:8}
(a) Elektronisch: $\gamma \approx \qty{5e-4}{J/(mol.K^2)}$ geeft
\qty{0.5}{mJ/(mol.K)}; rooster: $1944R\dots$ --- de $T^3$-term van
Debye geeft $\approx\qty{0.05}{mJ/(mol.K)}$. (b) De elektronen, met
een ordegrootte verschil. (c) Rond \qty{3}{K}
(\cref{exo:b3:quantum-statistics:5}). (d) Onder de omslag is het
gemeten snijpunt van $C/T$ gelijk aan $\gamma \propto g(E_{\text{F}})$:
een experiment met een thermometer en een verwarmingselement leest
de \omterm{prop:b3:schrodinger-three-dimensions:counting}{toestandsdichtheid} aan het fermioppervlak af.
\end{solution}

\begin{solution}{exo:b3:photons-phonons:9}
(a) Laag: absorbeert $\sigma T_s^4$, straalt $2\sigma T_a^4$ uit (beide
zijden): $T_s^4 = 2T_a^4$. Oppervlak: $S + \sigma T_a^4 = \sigma
T_s^4$. (b) Elimineer $T_a$: $S = \tfrac12\sigma T_s^4$, dus\
$T_s = 2^{1/4}T_{\text{eff}}$. (c) $1.19 \times 255 =
\qty{303}{K}$ --- ruim boven de werkelijke \qty{288}{K}, omdat de
echte deken gedeeltelijk en gelaagd is. (d) De vibratiekwanta
van drieatomige sporengassen: doorzichtig bij $0.5\,\unit{\micro m}$,
ondoorzichtig bij $15\,\unit{\micro m}$ ---
het molecuul van \cref{pb:b3:harmonic-oscillator:1}.
\end{solution}

\begin{solution}{exo:b3:photons-phonons:10}
(a) $u = (8\pi k_{\text{B}}^4T^4/h^3c^3)\int x^3/(\eu^x - 1)\dd x
= 8\pi^5k_{\text{B}}^4T^4/15h^3c^3$. (b--c) $\sigma = ac/4 =
2\pi^5k_{\text{B}}^4/15h^3c^2 = \qty{5.67e-8}{W.m^{-2}.K^{-4}}$
--- uit drie constanten, precies goed. (d) De \emph{vorm} van de kromme
legt $h/k_{\text{B}}$ vast (via $h\nu/k_{\text{B}}T$) en haar
absolute \emph{schaal} legt $k_{\text{B}}^4/h^3$ vast: twee
vergelijkingen, twee constanten --- Planck las $h$ \emph{en} $N_{\text{A}} =
R/k_{\text{B}}$ uit één fit af, jaren voor de korreltjes van Perrin.
\end{solution}

\begin{solution}{exo:b3:photons-phonons:11}
(a) $B_{12}uN_1 = B_{21}uN_2 + AN_2$. (b) Los $u$ op:
$u = (A/B_{21})/[(B_{12}/B_{21})\eu^{h\nu_0/k_{\text{B}}T} - 1]$;
overeenstemming met Plancks vorm voor alle $T$ dwingt $B_{12} = B_{21}$ en
$A/B = 8\pi h\nu_0^3/c^3$ af. (c) Zonder $A$ is er geen evenwicht
mogelijk tegen de \omterm{thm:b3:photons-phonons:planck}{wet van Planck} in: het evenwicht zelf eist dat
aangeslagen atomen ongevraagd vervallen, met een tempo dat aan het
gestimuleerde vastzit --- Einstein voorspelde spontane emissie uit
de thermodynamica, tien jaar voordat de kwantummechanica haar
afleidde. (d) $A \propto
\nu^3$: ultraviolette toestanden vervallen in nanoseconden
(\cref{exo:b3:perturbation-theory:8}), zodat inversies bij korte
golflengten ruïneus duur zijn --- het probleem van de
röntgenlaser.
\end{solution}

\begin{solution}{exo:b3:photons-phonons:12}
(a) $\kappa = \tfrac13 \times \num{1.8e6} \times \num{1.3e4}
\times \num{3e-7} \approx \qty{2300}{W/(m.K)}$: vijf keer
koper. (b) Stijf en licht betekent snelle fononen; een zuiver,
perfect, licht rooster biedt weinig om ze aan te verstrooien ---
$\ell$ haalt duizenden roosterconstanten. (c) Zelfs het
massaverschil van $^{13}$C-atomen verstrooit fononen meetbaar:
isotopenwanorde is een echte weerstand, die met zuivering
verdwijnt. (d) Warmtespreiders onder vermogenschips; en de warmte
van uw lip stroomt zo snel weg dat een echte diamant
verrassend koud aanvoelt --- de tastproef van de juwelier is een
fononmeting.
\end{solution}

\begin{solution}{pb:b3:photons-phonons:1}
\textbf{1.} $\lambda_{\max} = 2898/2.725 \approx \qty{1.06}{mm}$;
$411$ fotonen per cm$^3$.
\textbf{2.} $u = aT^4 = \qty{4.2e-14}{J/m^3} \approx
\qty{0.26}{eV/cm^3}$; massa-equivalent $u/c^2 \approx
\qty{4.6e-31}{kg/m^3}$: ongeveer $5\times10^{-5}$ van de kritieke waarde.
\textbf{3.} Een thermisch spectrum vereist dat de bron en de
straling in evenwicht zijn gekomen --- onmogelijk voor welke
verzameling sterren, stof of plasma's ook die langs
verschillende zichtlijnen wordt opgeteld; alleen een
heelal dat \emph{ooit zelf} een hete ondoorzichtige holte was, past.
\textbf{4.} $1 + z = 3000/2.725 \approx 1100$.
\textbf{5.} Saha's evenwicht sloeg om naar neutrale atomen: de vrije
elektronen verdwenen, de thomsonverstrooiing hield op, en de fotonen
vliegen sindsdien onaangeroerd verder.
\textbf{6.} Helder bij millimetergolflengten, bij
\qty{2.725}{K}: de hemel van Olbers \emph{staat} in lichterlaaie ---
de roodverschuiving heeft de gloed uit het zichtbare gehaald.
\textbf{7.} $\num{4.1e8}/0.25 \approx \num{1.6e9}$ fotonen per
baryon.
\textbf{8.} De fotonen zijn de as van de annihilatie van materie
en antimaterie; de baryonen zijn het overschot van één op een
miljard dat overleefde: $\eta$ meet de oorspronkelijke asymmetrie
--- nog altijd onverklaard.
\textbf{9.} $\qty{4.2e-14}{}$ tegenover $\qty{2e-15}{J/m^3}$: het
relictlicht overtreft in energiedichtheid alles wat de sterren
ooit hebben uitgezonden.
\textbf{10.} Straling verdunt als $a^{-4}$ (aantal $a^{-3}$,
elk foton uitgerekt met $a^{-1}$), materie als $a^{-3}$: vroeg
regeerde de straling; de materie nam het over bij de gelijkheid,
waarna structuur kon groeien.
\textbf{11.} $1 + z_{\text{eq}} \approx 3400$.
\textbf{12.} $T_\nu = (4/11)^{1/3} \times 2.725 =
\qty{1.95}{K}$: een voorspelde kosmische neutrinoachtergrond, de
stoutmoedigste niet-geïnde cheque van de thermische fysica.
\textbf{13.} $v = c\,\Delta T/T = \num{3e8} \times \num{1.25e-3}
\approx \qty{370}{km/s}$: de beweging van het zonnestelsel door
het relictstelsel.
\textbf{14.} De rimpelingen in temperatuur (en dus in dichtheid en
snelheid) van het heelal op de leeftijd van 380\,000 jaar: een
foto van de zaden.
\textbf{15.} Te dichte plekken trekken harder en groeien (de
gravitationele instabiliteit); vóór de gelijkheid veerde de
stralingsdruk ze terug --- de groei wachtte tot de materie regeerde.
\textbf{16.} De afstand die het geluid ($c/\sqrt3$ in de
foton--baryonvloeistof) had afgelegd bij de recombinatie --- de
geluidshorizon, een akoestische meetlat begraven in de hemel.
\textbf{17.} De hoek komt tot op ongeveer een procent overeen met
de meetkunde van een vlak heelal: evenwijdige lijnen blijven,
kosmisch gezien, evenwijdig.
\textbf{18.} Monopool: het hete begin. Dipool: onze
\qty{370}{km/s}. Rimpelingen: de oorsprong van structuur en de
meetkunde van de ruimte --- elke decimaal een ontdekking.
\textbf{19.} Hij was in elke richting gelijk, dag en nacht,
zomer en winter, duivenvrij: niets plaatselijks, galactisch of
instrumenteels is zo isotroop --- alleen de hemel zelf.
\textbf{20.} Een sissen op elk niet-afgestemd toestel: de oerknal
was al die tijd op televisie.
\textbf{21.} Hete clusterelektronen schoppen CMB-fotonen via
compton\-verstrooiing omhoog in frequentie: onder de omslag haalt
de cluster fotonen weg --- een silhouet waarvan de diepte niet van
de afstand afhangt, zodat surveys clusters bij elke
roodverschuiving vinden.
\textbf{22.} Oorspronkelijke zwaartekrachtgolven zouden de
polarisatie van de CMB tot ``B-moden'' verdraaien: de golven van de
relativiteitstheorie afgelezen in thomsonverstrooid licht --- twee
draden van dit boek in één knoop.
\textbf{23.} $n_\gamma c/4 \approx \qty{3e16}{m^{-2}s^{-1}}$:
ongeveer $10^{16}$ oerknalfotonen passeren u per seconde,
ongevoeld.
\textbf{24.} Het is een stelsel dat door de \emph{inhoud} wordt
uitgezonderd (waar de straling isotroop is), niet door de wetten:
de relativiteitstheorie verbiedt alleen stelsels die door de wetten
bevoorrecht worden --- bewegen door de CMB is waarneembaar,
bewegen door de ruimtetijd niet.
\textbf{25.} Een planckkromme van \qty{2.725}{K} op elke zichtlijn
--- $411$ fotonen per kubieke centimeter, $10^9$ per baryon,
rimpelingen op $10^{-5}$: het zwarte lichaam van de statistische
fysica, verheven tot het grondwettelijke document van de kosmologie.
\end{solution}
